\begin{theorem}[The height at the marked basepoint]
\hypertarget{tag-02QP}{}\label{night:P1.16}
\StatementRecord{night:P1.16}{proved}{02QP}{\SeriesRef{night:P1.15}, \SeriesRef{night:P1.18}, \SeriesRef{night:P1.8}}
In the canonical splitting $\sigma_b$ of \cref{night:P1.15},
\[
\begin{aligned}
-1.0342766224163217994942488453783461&<h_b(1),\\
h_b(1)&<-1.0342766224163217994942488453783460.
\end{aligned}
\]
In particular $h_b(1)<-1$. The parameter $c=5/(3\pi)$,
the exceptional value of \cref{night:P1A.2}, lies outside
the global transversality domain.
\end{theorem}
\begin{proof}
Use the root transports, normalized intertwiner and interval
linear system in \cref{night:P1.18}. At the marked basepoint
$H(1)=I$, so the real system in \cref{def:hb} uses $Y=X$.
Its enclosure has radius less than $7\cdot10^{-76}$ and is
contained in the following wider interval:
\[
\begin{aligned}
h_b(1)\in\bigl[&
-1.03427662241632179949424884537834605259496477790218821132\\
&\qquad\mathbin{\pm}8.11\cdot10^{-57}\bigr].
\end{aligned}
\]
This proves the two strict bounds. At zero the nontrivial deck
characters likewise force the normalized coefficients
$v_2,v_1$ to vanish. With the regular limiting frame $F_0$
of \cref{night:P1.18}, the same enclosure method gives
\[
\begin{aligned}
h_b(0)&=-\frac{96}{625\pi}\operatorname{Re}(F_0)_{03}\\
&=-0.931793376356523505430795105796011114990828\ldots .
\end{aligned}
\]

The connection of this calculation with based-loop continuation
is an exact gauge identity. Put $a=256/3125$ and
\[
D(\lambda)=\operatorname{diag}
 (1,\lambda^2,\lambda,1,\lambda^{-1},\lambda^{-2}).
\]
Then $\widetilde Y=YD$ satisfies the Fuchsian equation
\[
\begin{aligned}
\widetilde\Gamma
 &=D^{-1}\Gamma_6D+D^{-1}D'
 =\frac{R_0}{\lambda}+\frac{R_*}{\lambda+a},\\
R_0&=\lim_{\lambda\to0}\lambda\widetilde\Gamma,
\qquad
R_* =\lim_{\lambda\to-a}(\lambda+a)\widetilde\Gamma.
\end{aligned}
\]
Since $D(1)=I$, the based meridians of \cref{def:hb}
retain their matrices $B,A$. The regular-root deck calculation
in \cref{night:P1.18} represents those same meridians.
Both descriptions therefore give the same normalized
intertwiner $X$, by its uniqueness in \cref{def:hb}.
If $T_{12}$ is the Fuchsian transition along $[1,2]$ and
$u=2^{-1/5}$, uniqueness of the normalized differential equation
gives the exact equality
\[
XT_{12}D(2)^{-1}
 =F_\infty U_\infty(u)D_\infty(u)^{-1}.
\]
Thus the loop and regular-root calculations refer to the same
period plane with the same scalar and based marking.

Finally $5/(3\pi)<1$ and $h_b(1)<-1$, whereas
$h_b(\lambda)\to+\infty$ at the cusp. Along a path from $1$
to that region the function $h_b+5/(3\pi)$ has a zero.
The exact transversality criterion of \cref{night:P1.8}
therefore excludes this parameter.
\end{proof}
